\documentclass[../../main.tex]{subfiles}
\begin{document}

\tikzstyle{Node_Default}=[ draw, thick]
\tikzstyle{Edge_Default}=[-, thick]
\begin{problem}
    \begin{align*}
        n &= 4, \\
        (a_1, a_2, a_3, a_4) &= (do,if,stop,then), \\
        P(1:4) &=(3,3,3,1)\\
        Q(0:4) &=(1,3,2,1,1)\\
    \end{align*}
    计算\(W(i,j),R(i,j),C(i,j)\),并给出最优二分检索树.
\end{problem}
\begin{solution}
    \begin{align*}
        W(i,j) &= W(i,j-1) + P(j) + Q(j)\\
        C(i,j) &= \min_{i \leq k \leq j} \{C(i,k-1) + C(k,j) + W(i,j)\}\\
    \end{align*}
    算得表格如下:
    \begin{table}[H]
        \centering
        \begin{tabular}{cc|cccccc}
            &j-i& &  &  &  &  \\
            i&\(W,C,R\)& 0 &1 & 2 & 3 & 4  \\
            \hline
            &0&1,0,0&3,0,0&2,0,0&1,0,0&1,0,0\\
            &1&7,7,1&8,8,2&4,4,3&3,3,4& - \\
            &2&12,19,2&10,14,2&6,9,3& - & - \\
            &3&14,25,2&12,21,2& - & - & - \\
            &4&14,30,2& - & - & - & - \\
        \end{tabular}
    \end{table}
    最优二分检索树如下:
    \begin{center}
        \begin{tikzpicture}
            \node[Node_Default] (A) at ( 0 , 0 ) {if};
            \node[Node_Default] (B) at ( -1 , -1 ) {do};
            \node[Node_Default] (C) at ( 1 , -1 ) {stop};
            \node[Node_Default] (D) at ( 2 , -2 ) {then};
            \draw[Edge_Default] (A) -- (B);
            \draw[Edge_Default] (A) -- (C);
            \draw[Edge_Default] (C) -- (D);
            \end{tikzpicture}
    \end{center}
\end{solution}
\end{document}